\begin{trapbox}
\begin{description}[style=nextline, leftmargin=2.8em, labelsep=0.5em, itemsep=0.35em]
    \item[\textbf{\textcolor{CTrap}{易错点一（忽视非负条件）}}] 认为不等式 $\frac{a+b+c}{3} \ge \sqrt[3]{abc}$ 对所有实数 $a,b,c$ 都成立。
    \item[\textbf{\textcolor{CTrap}{正确理解（非负是前提）：}}] 该不等式成立的前提是 $a,b,c \ge 0$。若有负数，几何平均可能无意义（开三次方有意义但结论不成立），或者结论本身失效。
    \item[\textbf{\textcolor{CTrap}{错因分析（混淆适用范围）：}}] 忽略了算术平均数对负数也成立，但几何平均数要求非负。不满足非负条件时使用，会导致推导错误。
\end{description}

\medskip
\begin{description}[style=nextline, leftmargin=2.8em, labelsep=0.5em, itemsep=0.35em]
    \item[\textbf{\textcolor{CTrap}{易错点二（乱用定值条件）}}] 看到 $a+b+c$ 或 $abc$ 就直接套公式求最值，没有验证“和”或“积”是否为定值。
    \item[\textbf{\textcolor{CTrap}{正确理解（定值是目标）：}}] 只有在“和为定值”时，才能用此式求 $abc$ 的最大值；在“积为定值”时，才能求 $a+b+c$ 的最小值。否则不等式本身成立，但无法确定最值。
    \item[\textbf{\textcolor{CTrap}{错因分析（误解应用逻辑）：}}] 混淆了“不等式恒成立”和“基于不等式求最值”两个不同问题。求最值时，必须存在一个在取等条件下能达到的定值边界。
\end{description}

\medskip
\begin{description}[style=nextline, leftmargin=2.8em, labelsep=0.5em, itemsep=0.35em]
    \item[\textbf{\textcolor{CTrap}{易错点三（忽略取等验证）}}] 求出最值后，不验证等号成立条件是否可以被满足，或默认存在。
    \item[\textbf{\textcolor{CTrap}{正确理解（取等须可达）：}}] 用不等式 $A \ge B$ 求最小值 $B$ 时，必须检查是否存在一组变量使得 $A=B$ 成立。否则 $B$ 只是下界而非最小值（最小值可能更大）。
    \item[\textbf{\textcolor{CTrap}{错因分析（逻辑不严谨）：}}] 误认为由不等式推导出的边界一定能够取到。如果取等条件与题目其他约束矛盾，则最值在别处取得，需改用其他方法。
\end{description}
\end{trapbox}
